\chapter{Titik Tengah dan Garis Sejajar}\label{ch:g8:midpoints}

Hubungkan titik tengah dua sisi sebuah segitiga: \omterm{def:g6:lines:objects}{ruas garis} yang kamu
peroleh selalu \omterm{def:g6:lines:perp}{sejajar} dengan sisi yang ketiga, dan panjangnya tepat
setengahnya. ``Teorema titik tengah'' ini beserta konversnya adalah
cicipan pertama sebuah gagasan besar --- \omterm{def:g6:lines:perp}{garis sejajar} memotong
panjang secara \omterm{def:g6:propdata:def}{sebanding} --- yang mekar menjadi teorema Thales pada
\cref{ch:g9:thales}.

\section{Teorema titik tengah}

\begin{theorem}[Teorema titik tengah]\label{thm:g8:midpoints:direct}
Pada segitiga $ABC$, misalkan $I$ titik tengah $[AB]$ dan $J$ titik
tengah $[AC]$. Maka garis $(IJ)$ \omterm{def:g6:lines:perp}{sejajar} dengan $(BC)$, dan
\[
IJ = \frac{BC}{2} .
\]
\end{theorem}

\begin{proof}[Gagasan buktinya]
Misalkan $K$ titik \omterm{def:g7:central:def}{simetris} $I$ terhadap $J$ (yaitu setengah putaran,
\cref{ch:g7:central}). \omterm{def:g4:geometry:polygon}{Segi empat} $AICK$ punya diagonal $[IK]$
dan $[AC]$ yang berpotongan di titik tengah bersamanya $J$: jadi ia
\omterm{def:g7:central:parallelogram}{jajargenjang}, sehingga $KC$ \omterm{def:g6:lines:perp}{sejajar} dengan $AI$, yaitu dengan $(AB)$,
dan $KC = AI = IB$. Lalu $IBCK$ punya dua sisi berhadapan ($[IB]$ dan
$[KC]$) yang \omterm{def:g6:lines:perp}{sejajar} dan sama panjang: jadi ia pun \omterm{def:g7:central:parallelogram}{jajargenjang}
(\cref{thm:g7:central:recognize}), sehingga $(IK) \parallel (BC)$ ---
yaitu $(IJ) \parallel (BC)$ --- dan $BC = IK = 2\,IJ$.
\end{proof}

\begin{omfigure}
\begin{tikzpicture}[scale=1.1]
  \coordinate (A) at (1.7,2.9);
  \coordinate (B) at (0,0);
  \coordinate (C) at (4.6,0);
  \coordinate (I) at ($(A)!0.5!(B)$);
  \coordinate (J) at ($(A)!0.5!(C)$);
  \draw[thick] (A) -- (B) -- (C) -- cycle;
  \draw[omThm, very thick] (I) -- (J);
  \fill (A) circle (1.5pt) node[above, font=\small] {$A$};
  \fill (B) circle (1.5pt) node[below left, font=\small] {$B$};
  \fill (C) circle (1.5pt) node[below right, font=\small] {$C$};
  \fill[omThm] (I) circle (1.5pt) node[left, font=\small] {$I$};
  \fill[omThm] (J) circle (1.5pt) node[right, font=\small] {$J$};
  \draw ($(A)!0.22!(B)$) ++(-0.09,-0.05) -- ++(0.18,0.1);
  \draw ($(A)!0.72!(B)$) ++(-0.09,-0.05) -- ++(0.18,0.1);
  \draw ($(A)!0.25!(C)$) ++(-0.06,-0.1) -- ++(0.12,0.2);
  \draw ($(A)!0.25!(C)$) ++(0.02,-0.1) -- ++(0.12,0.2);
  \draw ($(A)!0.75!(C)$) ++(-0.06,-0.1) -- ++(0.12,0.2);
  \draw ($(A)!0.75!(C)$) ++(0.02,-0.1) -- ++(0.12,0.2);
\end{tikzpicture}

{\small \omterm{def:g6:lines:objects}{Ruas garis} titik tengah $[IJ]$ (merah) \omterm{def:g6:lines:perp}{sejajar} dengan sisi
ketiganya $(BC)$ dan panjangnya setengahnya.}
\end{omfigure}

\begin{example}\label{ex:g8:midpoints:direct}
Pada segitiga $ABC$ dengan $BC = 9$ cm, titik tengah $I$ dari $[AB]$
dan $J$ dari $[AC]$ dihubungkan. Tanpa pengukuran apa pun:
$(IJ) \parallel (BC)$ dan $IJ = \frac92 = 4.5$ cm. Segitiga kecil
$AIJ$ adalah salinan $ABC$ berukuran setengahnya --- \omterm{def:g6:measure:perimeter}{kelilingnya} pun
setengah \omterm{def:g6:measure:perimeter}{keliling} $ABC$.
\end{example}

\begin{theorem}[Konvers]\label{thm:g8:midpoints:converse}
Pada segitiga $ABC$, garis yang melalui titik tengah $I$ dari $[AB]$
dan \emph{\omterm{def:g6:lines:perp}{sejajar}} dengan $(BC)$ memotong sisi $[AC]$ di titik tengahnya.
\end{theorem}
\admitted

\begin{example}\label{ex:g8:midpoints:converse}
Sebuah jalan setapak menyeberangi ladang berbentuk segitiga, bermula
di tengah salah satu tepinya lalu berjalan \omterm{def:g6:lines:perp}{sejajar} dengan tepi yang di
seberangnya. Konversnya menjamin bahwa jalannya keluar tepat di tengah
tepi yang lain --- tanpa perlu mengukur di sisi yang jauh.
\end{example}

\begin{method}[Memilih antara teorema dan konversnya]\label{met:g8:midpoints:choose}
\begin{enumerate}
  \item \emph{Dua titik tengah diketahui} $\Rightarrow$ pakailah
        teorema langsungnya: simpulkan kesejajaran dan setengah
        panjangnya.
  \item \emph{Satu titik tengah dan satu \omterm{def:g6:lines:perp}{garis sejajar} diketahui}
        $\Rightarrow$ pakailah konversnya: simpulkan bahwa titik
        potongnya adalah titik tengah.
  \item Tulislah segitiga yang mana, titik tengah yang mana dan
        teorema yang mana yang kamu pakai --- satu kalimat untuk
        masing-masing.
\end{enumerate}
\end{method}

\section{Segitiga berukuran setengah}

\begin{proposition}[Segitiga titik tengah]\label{prop:g8:midpoints:medial}
Menghubungkan ketiga titik tengah sisi sebuah segitiga memotongnya
menjadi empat segitiga yang luasnya sama, masing-masing salinan
aslinya berukuran setengah.
\end{proposition}

\begin{proof}
Tiap sisi segitiga titik tengahnya \omterm{def:g6:lines:perp}{sejajar} dengan sebuah sisi $ABC$
dan panjangnya setengahnya (\cref{thm:g8:midpoints:direct}, yang
dipakai tiga kali). Keempat segitiga kecilnya punya sisi dengan ketiga
panjang yang sama, jadi keempatnya salinan yang identik; bersama-sama
keempatnya mengubin $ABC$, jadi masing-masing punya \omterm{def:g3:division:half}{seperempat}
luasnya.
\end{proof}

\begin{omfigure}
\begin{tikzpicture}[scale=1.05]
  \coordinate (A) at (1.9,3.0);
  \coordinate (B) at (0,0);
  \coordinate (C) at (5.0,0);
  \coordinate (I) at ($(A)!0.5!(B)$);
  \coordinate (J) at ($(A)!0.5!(C)$);
  \coordinate (K) at ($(B)!0.5!(C)$);
  \draw[thick] (A) -- (B) -- (C) -- cycle;
  \fill[omThm!12] (I) -- (J) -- (K) -- cycle;
  \draw[omThm, thick] (I) -- (J) -- (K) -- cycle;
  \fill (A) circle (1.4pt) node[above, font=\small] {$A$};
  \fill (B) circle (1.4pt) node[below left, font=\small] {$B$};
  \fill (C) circle (1.4pt) node[below right, font=\small] {$C$};
  \fill[omThm] (I) circle (1.4pt) node[left, font=\small] {$I$};
  \fill[omThm] (J) circle (1.4pt) node[right, font=\small] {$J$};
  \fill[omThm] (K) circle (1.4pt) node[below, font=\small] {$K$};
\end{tikzpicture}

{\small Segitiga titik tengah $IJK$ memotong $ABC$ menjadi empat
segitiga yang sama --- gambar yang layak diingat.}
\end{omfigure}

\begin{remark}[Menuju Thales]
Teorema titik tengah adalah keadaan ``satu setengah'' dari kenyataan
yang lebih umum: garis yang \omterm{def:g6:lines:perp}{sejajar} dengan salah satu sisi segitiga
memotong kedua sisi lainnya dengan \emph{perbandingan yang sama} ---
sepertiga dan sepertiga, dua perlima dan dua
perlima\,\dots\ Pernyataan umum itu adalah teorema Thales
(\cref{ch:g9:thales}); keadaan titik tengahnya satu-satunya yang
diperlukan tahun ini.
\end{remark}

\section{Latihan}

\begin{exercise}[$\star$]\label{exo:g8:midpoints:1}
Pada segitiga $RST$, $M$ adalah titik tengah $[RS]$ dan $N$ titik
tengah $[RT]$, dengan $ST = 7$ cm. Apa yang dapat dikatakan tentang
garis $(MN)$ dan panjang $MN$? Sebutkan teorema yang kamu pakai.
\end{exercise}

\begin{exercise}[$\star$]\label{exo:g8:midpoints:2}
Pada segitiga $ABC$, $I$ adalah titik tengah $[AB]$ dan garis melalui
$I$ yang \omterm{def:g6:lines:perp}{sejajar} dengan $(BC)$ memotong $[AC]$ di $J$, dengan
$AC = 11$ cm. Hitunglah $AJ$, sambil menyebut teorema yang dipakai.
\end{exercise}

\begin{exercise}[$\star$]\label{exo:g8:midpoints:3}
Gambarlah segitiga $ABC$ dengan $AB = 8$ cm, $AC = 6$ cm, $BC = 7$ cm,
letakkan titik tengah $I$ dari $[AB]$ dan $J$ dari $[AC]$, lalu ukurlah
$IJ$ untuk memeriksa teoremanya. Hitunglah \omterm{def:g6:measure:perimeter}{keliling} $AIJ$ tanpa
mengukur apa pun lagi.
\end{exercise}

\begin{exercise}[$\star$]\label{exo:g8:midpoints:4}
$IJK$ adalah segitiga titik tengah $ABC$, yang sisinya berukuran $10$,
$12$ dan $16$ cm. Sebutkan ketiga sisi $IJK$ dan \omterm{def:g6:measure:perimeter}{kelilingnya}.
Bagaimana perbandingan kedua \omterm{def:g6:measure:perimeter}{kelilingnya}?
\end{exercise}

\begin{exercise}[$\star$]\label{exo:g8:midpoints:5}
Luas sebuah segitiga adalah $36$ cm$^2$. Berapa luas segitiga titik
tengahnya? Berapa luas tiap segitiga kecil dari keempatnya?
\end{exercise}

\begin{exercise}[$\star\star$]\label{exo:g8:midpoints:6}
Pada segitiga $ABC$, $I$ adalah titik tengah $[AB]$, $J$ titik tengah
$[AC]$ dan $K$ titik tengah $[BC]$. Tunjukkan bahwa $IJKB$ adalah
\omterm{def:g7:central:parallelogram}{jajargenjang}. (Bandingkan $[IJ]$ dan $[BK]$: sejajarkah? samakah?)
\end{exercise}

\begin{exercise}[$\star\star$]\label{exo:g8:midpoints:7}
Sebuah layar berbentuk segitiga punya tepi terbawah sepanjang $6.4$ m.
Sebuah jahitan penguat menghubungkan titik tengah kedua tepi lainnya.
Berapa panjang jahitan yang diperlukan? Bagaimana kalau jahitannya
justru menghubungkan titik yang terletak pada \omterm{def:g3:division:half}{seperempat} tiap tepinya,
bermula dari \omterm{def:g5:solids:def}{titik sudut} yang di atas? (Dugalah dengan gambar;
buktinya milik Thales, \cref{ch:g9:thales}.)
\end{exercise}

\begin{exercise}[$\star\star$]\label{exo:g8:midpoints:8}
Pada segitiga $ABC$ yang siku-siku di $A$, $M$ adalah titik tengah
sisi miringnya $[BC]$, dan \omterm{def:g6:lines:perp}{garis sejajar} $(AB)$ yang melalui $M$
memotong $[AC]$ di $N$.
\begin{enumerate}
  \item Tunjukkan bahwa $N$ adalah titik tengah $[AC]$.
  \item Jelaskan mengapa $(MN)$ \omterm{def:g6:lines:perp}{tegak lurus} terhadap $(AC)$.
\end{enumerate}
\end{exercise}

\begin{exercise}[$\star\star\star$]\label{exo:g8:midpoints:9}
Misalkan $ABCD$ \omterm{def:g4:geometry:polygon}{segi empat} \emph{apa pun}, dan $I$, $J$, $K$, $L$
titik tengah $[AB]$, $[BC]$, $[CD]$, $[DA]$ (yaitu dugaan pada
\cref{exo:g7:central:11}).
\begin{enumerate}
  \item Terapkan teorema titik tengah pada segitiga $ABC$ untuk \omterm{def:g6:lines:objects}{ruas
        garis} $[IJ]$, dan pada segitiga $ACD$ untuk \omterm{def:g6:lines:objects}{ruas garis}
        $[LK]$: bandingkan masing-masing dengan diagonal $[AC]$.
  \item Simpulkan bahwa $IJKL$ selalu \omterm{def:g7:central:parallelogram}{jajargenjang}.
\end{enumerate}
\end{exercise}

\section{Soal: Ketiga garis berat dan titik berat}

\begin{problem}[{Soal akhir pekan --- garis berat sebuah segitiga
bertemu di satu titik, pada dua pertiga sepanjang masing-masingnya}]
\label{pb:g8:midpoints:1}
Potonglah sebuah segitiga dari karton yang kaku dan ia akan
seimbang, benar-benar datar, di atas ujung pensil yang diletakkan
pada satu titik istimewa: \emph{titik berat}\index{titik berat}
segitiganya. Soal ini mencari titik itu. \emph{Garis berat}\index{garis
berat} sebuah segitiga adalah \omterm{def:g6:lines:objects}{ruas garis} yang menghubungkan \omterm{def:g5:solids:def}{titik
sudut} dengan \emph{titik tengah} sisi yang di seberangnya. Kamu akan
membuktikan --- dengan teorema titik tengah yang dipakai dua kali,
dengan cara yang tidak terduga --- bahwa ketiga \omterm{pb:g8:midpoints:1}{garis beratnya}
melewati satu titik, lalu menemukan letaknya dengan tepat.

\medskip
\noindent\textbf{Bagian I --- Pemanasan.}
\begin{enumerate}
  \item Gambarlah segitiga $ABC$ yang besar (tanpa bentuk khusus),
        lukislah titik tengah $I$ dari $[AB]$, $J$ dari $[AC]$ dan $K$
        dari $[BC]$, lalu gambarlah ketiga \omterm{pb:g8:midpoints:1}{garis beratnya} $[AK]$,
        $[BJ]$, $[CI]$. Apa yang kamu amati?
  \item Buktikan bahwa sebuah \omterm{pb:g8:midpoints:1}{garis berat} memotong segitiganya
        menjadi dua segitiga yang luasnya sama. (Bandingkan alas dan
        tingginya, lalu pakailah rumus luasnya: alas $\times$ tinggi
        $\div\ 2$.) Untuk segitiga seluas $36$ cm$^2$, berapa luas
        kedua kepingnya?
  \item Pada segitiga $ABC$, apa yang dikatakan
        \cref{thm:g8:midpoints:direct} tentang \omterm{def:g6:lines:objects}{ruas garis}
        $[IJ]$?
\end{enumerate}

\medskip
\noindent\textbf{Bagian II --- Dua \omterm{pb:g8:midpoints:1}{garis berat} bertemu pada dua
pertiganya.}
\omterm{pb:g8:midpoints:1}{Garis berat} $[BJ]$ dan $[CI]$ berpotongan di sebuah titik; sebutlah
$G$. Misalkan $M$ titik tengah $[GB]$ dan $N$ titik tengah $[GC]$.
\begin{enumerate}[resume]
  \item Terapkan teorema titik tengah \emph{pada segitiga $GBC$}:
        apa yang dikatakannya tentang \omterm{def:g6:lines:objects}{ruas garis} $[MN]$?
  \item Simpulkan bahwa $(IJ) \parallel (MN)$ dan $IJ = MN$, lalu
        simpulkan dengan \cref{thm:g7:central:recognize} bahwa $IJNM$
        adalah \omterm{def:g7:central:parallelogram}{jajargenjang}.
  \item Jelaskan mengapa $I$ dan $N$ keduanya terletak pada \omterm{pb:g8:midpoints:1}{garis
        berat} $(CI)$, dan $J$ dan $M$ keduanya pada \omterm{pb:g8:midpoints:1}{garis berat}
        $(BJ)$ --- jadi diagonal \omterm{def:g7:central:parallelogram}{jajargenjang} $IJNM$ adalah $[IN]$
        dan $[JM]$, dan keduanya berpotongan tepat di $G$. Lalu apa
        yang dikatakan \emph{definisi} \omterm{def:g7:central:parallelogram}{jajargenjang} tentang $G$?
  \item Simpulkan bahwa $BM = MG = GJ$, lalu simpulkan:
        \[
        BG = 2 \times GJ
        \]
        --- titik $G$ duduk pada \omterm{pb:g8:midpoints:1}{garis berat} $[BJ]$ pada dua
        pertiga jalan dari \omterm{def:g5:solids:def}{titik sudut} $B$. Nyatakan dan berilah
        alasan untuk hasil serupa pada \omterm{pb:g8:midpoints:1}{garis berat} $[CI]$.
  \item Sekarang lupakan $[CI]$ lalu jalankan alasan yang persis sama
        pada pasangan \omterm{pb:g8:midpoints:1}{garis berat} $[BJ]$ dan $[AK]$: titik potongnya
        juga duduk pada dua pertiga $[BJ]$ dari $B$. Jelaskan mengapa
        itu memaksanya menjadi titik $G$ yang \emph{sama} --- lalu
        simpulkan bahwa ketiga \omterm{pb:g8:midpoints:1}{garis beratnya} melewati $G$.
  \item Pada segitiga yang \omterm{pb:g8:midpoints:1}{garis beratnya} $[BJ]$ berukuran $9$ cm,
        berapa jarak $G$ dari $B$, dan dari $J$?
  \item Berilah alasan (satu kalimat) bahwa $AG = 2 \times GK$ juga.
\end{enumerate}

\medskip
\noindent\textbf{Bagian III --- Enam irisan yang sama, dan koordinat.}
Ketiga \omterm{pb:g8:midpoints:1}{garis berat} itu memotong segitiganya menjadi enam segitiga
kecil di sekeliling $G$.
\begin{enumerate}[resume]
  \item Pada segitiga $GBC$, \omterm{def:g6:lines:objects}{ruas garis} $[GK]$ adalah \omterm{pb:g8:midpoints:1}{garis berat}.
        Simpulkan bahwa kedua segitiga kecil $GBK$ dan $GKC$ punya
        luas yang sama.
  \item Dengan memakai pertanyaan 2 pada segitiga $ABK$ dan $AKC$
        (keduanya dipotong \omterm{pb:g8:midpoints:1}{garis berat} $[AK]$ milik $ABC$),
        tunjukkan bahwa segitiga $ABG$ dan $ACG$ punya luas yang
        sama. Dengan mengulanginya memakai \omterm{pb:g8:midpoints:1}{garis berat} yang lain,
        simpulkan bahwa ketiga segitiga
        $ABG$, $BCG$, $CAG$ semuanya berluas sepertiga luas $ABC$.
  \item Masing-masing dari ketiga segitiga itu terpotong menjadi dua
        oleh sepotong \omterm{pb:g8:midpoints:1}{garis berat} (misalnya $[GI]$ adalah \omterm{pb:g8:midpoints:1}{garis berat}
        segitiga $ABG$ --- dilihat dari \omterm{def:g5:solids:def}{titik sudut} yang mana?).
        Simpulkan: keenam irisan kecil di sekeliling $G$ semuanya
        punya luas yang sama, yaitu seperenam segitiganya.
  \item Pada kisi berkoordinat, gambarlah $A(0, 0)$, $B(6, 0)$ dan
        $C(0, 6)$, dengan $K(3, 3)$ titik tengah $[BC]$. Pakailah
        teorema dua pertiga pada Bagian II untuk menghitung koordinat
        $G$ pada \omterm{pb:g8:midpoints:1}{garis berat} $[AK]$; lalu periksalah bahwa titik yang
        sama duduk pada dua pertiga \omterm{pb:g8:midpoints:1}{garis berat} $[BJ]$, dengan
        $J(0, 3)$. Bandingkan koordinat $G$ dengan rata-rata
        \[
        \frac{0 + 6 + 0}{3},
        \qquad
        \frac{0 + 0 + 6}{3} .
        \]
  \item Penutup tentang keseimbangannya: dengan memakai pertanyaan 2
        dan 8, jelaskan mengapa segitiga kartonnya seimbang di atas
        mata pisau yang dibaringkan sepanjang \omterm{pb:g8:midpoints:1}{garis berat} mana pun
        --- karton yang sama banyak pada tiap sisinya --- dan
        mengapa karena itu ujung pensilnya harus diletakkan di $G$,
        yaitu titik yang disebut para fisikawan sebagai \omterm{pb:g8:midpoints:1}{titik berat}.
        (Bukti keseimbangan yang benar-benar ketat perlu hitung
        integral pada volume di tingkat universitas; luas yang sama
        membuatnya masuk akal untuk hari ini.)

\end{enumerate}
\end{problem}
